% !TeX root = main.tex
\section{Aerodynamic and Simulation Model}
\subsection{Quasi-Steady Strip Theory}
We use a quasi-steady strip-theory aerodynamic model inspired by DeLaurier-style formulations, discretizing the wing span into strips and summing circulatory and added-mass contributions to obtain instantaneous force and moment.
The model also includes chordwise contributions (leading-edge suction, camber-related drag, and skin-friction drag) to predict the streamwise force component.
For computational efficiency and reproducibility, we emphasize cycle-averaged metrics computed from the simulated time series.

\subsection{Separation Modeling and Scope}
The DeLaurier-style model admits an optional separated-flow approximation triggered by a leading-edge dynamic-stall criterion.
In practice, the stall thresholds and separation parameters are highly uncertain for membrane wings and for the present kinematics.
As a result, enabling separation can qualitatively affect the existence of trim-like operating points.
In the main experiments of this paper, we disable the separated-flow model to obtain a controlled baseline and a consistent trim-like condition; separation sensitivity is discussed as a limitation.

\subsection{Wind-Tunnel Rig and Coordinate Convention}
We evaluate the model in an IsaacLab wind-tunnel rig where the body is fixed and the wings are driven by prescribed kinematics.
Aerodynamic forces are reported in the world frame, with the streamwise thrust direction defined as world $+X$ and lift defined as world $+Z$.
The freestream wind is aligned with world $-X$.

\subsection{Prescribed Twist}
To capture a basic coupling between stroke rate and spanwise twist, we use a prescribed twist model that scales the wing-tip twist magnitude with the commanded stroke rate.
The tip twist is clamped to a maximum magnitude and applied with a simple spanwise shape.
This serves as a tractable proxy for passive twist effects without requiring a full aeroelastic simulation.
